% ABA specification, Transition System 1 rule graph (PLTS.ABA.SpecStep).
% Four control states keyed on (mode, val). Rule 3 is drawn as a
% probabilistic branch: the undirected stem carries the tau-step, and the
% three outcome edges carry the flip masses. Five of the eight rules are
% drawn; fail (D1) and the corrupted interface callByz/retByz (D23) are not,
% none of them writing the mode. See the caption.
\begin{tikzpicture}[
    ->, >={Stealth[round]}, auto, semithick, on grid,
    every state/.style={draw, minimum size=1.15cm, align=center, font=\small},
    branch/.style={circle, draw, fill, inner sep=1.6pt},
    lbl/.style={font=\footnotesize, align=center, fill=white, inner sep=1.5pt},
    plbl/.style={font=\footnotesize, align=center, fill=white, inner sep=1.5pt}]

  \node[state, initial, initial text={}] (idl) at (0,0)      {idle\\$val{=}\bot$};
  \node[branch]                          (br)  at (3.9,0)    {};
  \node[state]                           (dea) at (7.4,2.0)  {terminal};
  \node[state]                           (loc) at (7.4,-2.0) {locked};
  \node[state, accepting]                (dec) at (3.9,-4.5) {decided\\$val$ set};

  % self-loops
  \path
    (idl) edge[loop above, looseness=5]
      node[lbl, above]{1 \texttt{callSet}\\2 \texttt{callLoop}} (idl)
    (dec) edge[loop below, looseness=5]
      node[lbl, below]{5 \texttt{ret}} (dec);

  % the flip as a probabilistic branch
  \draw[-] (idl) -- node[lbl, above, pos=0.45]{3 \texttt{coinFlip}, $\tau$} (br);
  \path
    (br) edge[bend right=14] node[plbl, left]{$\delta_{\mathrm f}$ (undelivered)} (dea)
    (br) edge[bend left=14]  node[plbl, right]{$\varepsilon$ (lock)} (loc)
    (br) edge[bend right=55, looseness=1.3]
      node[plbl, above]{$1{-}\varepsilon{-}\delta_{\mathrm f}$ (release)} (idl);

  % the decisions
  \path
    (loc) edge[bend left=12]  node[lbl, below]{4 \texttt{decide}} (dec);
  \draw[->] (idl) to[bend right=28]
    node[lbl, below]{4 \texttt{decide}} (dec);
\end{tikzpicture}
